% Part: turing-machines
% Chapter: machines-computations
% Section: configuration

\documentclass[../../../include/open-logic-section]{subfiles}

\begin{document}

\olfileid{tur}{mac}{con}
\olsection{التهيئات والحسابات}

\begin{explain}
تذكر كيف تتبعنا سابقًا تهيئات آلة الزوجية. إن الآلية المتخيلة المؤلفة
من الشريط ورأس القراءة والكتابة وبرنامج آلة تورنغ ليست في الحقيقة إلا
طريقة حدسية لتصور ماهية حساب آلة تورنغ. ويمكننا صوريًا تعريف حساب آلة
تورنغ عند دخل معطى بوصفه متتالية من \emph{التهيئات}؛ والتهيئة نفسها
ثلاثية تتألف من متتالية من الرموز (تقابل محتويات الشريط عند نقطة معينة
من الحساب)، وعدد يدل على موضع رأس القراءة والكتابة، وحالة. وباستعمال هذه العناصر
يمكننا تعريف ما تحسبه آلة تورنغ~$M$ عند دخل معطى.
\end{explain}

\begin{defn}[التهيئة]
\emph{تهيئة} آلة تورنغ $M=\tuple{Q,\Sigma,q_0,\delta}$ هي ثلاثية
$\tuple{C,m,q}$ حيث
\begin{enumerate}
\item $C \in \Sigma^*$ متتالية منتهية من رموز $\Sigma$،
\item $m \in \Nat$ عدد أصغر من $\len{C}$، و
\item $q \in Q$.
\end{enumerate}
وحدسيًا، المتتالية~$C$ هي محتوى الشريط (رموز جميع الخانات من الخانة
الواقعة في أقصى اليسار إلى آخر خانة غير فارغة أو سبق أن زارها الرأس)،
و$m$ رقم الخانة التي يمسحها رأس القراءة والكتابة (ابتداءً من ترقيم
الخانة الواقعة في أقصى اليسار بالعدد $0$)، و$q$ حالة الآلة الحالية.
\end{defn}

\begin{explain}
الدخل الممكن لآلة تورنغ متتالية من الرموز، وعادة ما تكون متتالية ترمز
عددًا على نحو ما. والتهيئة الابتدائية لآلة تورنغ هي التهيئة التي نبدأ
فيها تشغيل الآلة على ذلك الدخل: يحتوي الشريط علامة نهاية الشريط،
ويتبعها مباشرة الدخل المكتوب في الخانات الواقعة إلى يمينها؛ ويمسح
رأس القراءة والكتابة الخانة الواقعة في أقصى يسار الدخل (أي الخانة
الواقعة إلى يمين علامة الطرف الأيسر)؛ وتكون الآلية في الحالة
الابتدائية المعينة~$q_0$.
\end{explain}

\begin{defn}[التهيئة الابتدائية]
\emph{التهيئة الابتدائية} للآلة $M$ عند الدخل $I \in \Sigma^*$ هي
\[
\tuple{\TMendtape \frown I \frown \TMblank, 1, q_0}.
\]
\end{defn}

\begin{explain}
الرمز~$\frown$ يدل على \emph{الوصل}؛ وتبدأ سلسلة الدخل مباشرة إلى
يمين علامة نهاية الشريط.
\end{explain}

\begin{defn}
نقول إن التهيئة $\tuple{C,m,q}$ \emph{تنتج التهيئة
$\tuple{C',m',q'}$ في خطوة واحدة} (وفق الآلة~$M$) إذا وفقط إذا
\begin{enumerate}
\item كان الرمز رقم $m$ في $C$ هو $\sigma$،
\item عينت مجموعة تعليمات $M$ أن $\delta(q,\sigma)=
  \tuple{q',\sigma',D}$،
\item كان الرمز رقم $m$ في $C'$ هو $\sigma'$، و
\item
\begin{enumerate}
\item كان $D=L$ و$m'=m-1$ إذا كان $m>0$، وإلا كان $m'=0$، أو
\item كان $D=R$ و$m'=m+1$، أو
\item كان $D=N$ و$m'=m$،
\end{enumerate}
\item إذا كان $m'=\len{C}$، فإن $\len{C'}=\len{C}+1$ والرمز رقم
  $m'$ في $C'$ هو~$\TMblank$. وإلا فإن $\len{C'}=\len{C}$.
\item لكل $i$ بحيث $i<\len{C}$ و$i\neq m$، يكون $C'(i)=C(i)$.
\end{enumerate}
\end{defn}

\begin{defn}\ollabel{defn:run-output}
\emph{تشغيل $M$ عند الدخل~$I$} متتالية منتهية أو غير منتهية $C_i$
من تهيئات $M$، حيث $C_0$ هي التهيئة الابتدائية للآلة $M$ عند
الدخل~$I$، وتنتج كل $C_i$ التهيئة $C_{i+1}$ في خطوة واحدة متى وجدت
$C_{i+1}$. وإذا كانت المتتالية منتهية، فلا تنتج تهيئتها الأخيرة أي
تهيئة في خطوة واحدة.

نقول إن $M$ \emph{تتوقف عند الدخل $I$ بعد $k$ خطوات} إذا كان التشغيل
منتهيًا عند $C_k=\tuple{C,m,q}$، وكان الرمز رقم $m$ في~$C$ هو
$\sigma$، وكانت $\delta(q,\sigma)$ غير معرّفة. وفي هذه الحال يكون
\emph{خرج}~$M$ عند الدخل~$I$ هو~$O$، حيث $O$ سلسلة رموز لا تنتهي
بـ~$\TMblank$ بحيث $C=\TMendtape\concat O\concat\TMblank^j$ لبعض
$j \in \Nat$. (و$\TMblank^j$ متتالية من $j$ من رموز $\TMblank$.)
\end{defn}

\begin{explain}
وفق هذا التعريف ينتهي خرج~$O$ للآلة~$M$ دائمًا برمز غير~$\TMblank$؛
أو يكون $O$ السلسلة الخالية إذا كان الشريط كله عند الزمن~$k$ مملوءًا
برموز~$\TMblank$ (باستثناء $\TMendtape$ الواقع في أقصى اليسار).
\end{explain}

\end{document}
